\documentclass[10pt,a4paper]{amsart}
\usepackage{amsmath,amssymb,amsthm,graphicx}
\usepackage{verbatim}
\newenvironment{reference}{\comment}{\endcomment}
%\newenvironment{reference}{}{}
\newenvironment{slogan}{\comment}{\endcomment}
\newenvironment{history}{\comment}{\endcomment}

% For commutative diagrams we use Xy-pic
\usepackage[all]{xy}

% We use 2cell for 2-commutative diagrams.
\xyoption{2cell}
% Theorem environments.
%
\theoremstyle{plain}
\newtheorem{theorem}[subsection]{Theorem}
\newtheorem{proposition}[subsection]{Proposition}
\newtheorem{lemma}[subsection]{Lemma}

\theoremstyle{definition}
\newtheorem{definition}[subsection]{Definition}
\newtheorem{example}[subsection]{Example}
\newtheorem{exercise}[subsection]{Exercise}
\newtheorem{situation}[subsection]{Situation}

\theoremstyle{remark}
\newtheorem{remark}[subsection]{Remark}
\newtheorem{remarks}[subsection]{Remarks}

\numberwithin{equation}{subsection}

% Macros
%
\def\lim{\mathop{\mathrm{lim}}\nolimits}
\def\colim{\mathop{\mathrm{colim}}\nolimits}
\def\Spec{\mathop{\mathrm{Spec}}}
\def\Hom{\mathop{\mathrm{Hom}}\nolimits}
\def\Ext{\mathop{\mathrm{Ext}}\nolimits}
\def\SheafHom{\mathop{\mathcal{H}\!\mathit{om}}\nolimits}
\def\SheafExt{\mathop{\mathcal{E}\!\mathit{xt}}\nolimits}
\def\Sch{\mathit{Sch}}
\def\Mor{\mathop{\mathrm{Mor}}\nolimits}
\def\Ob{\mathop{\mathrm{Ob}}\nolimits}
\def\Sh{\mathop{\mathit{Sh}}\nolimits}
\def\NL{\mathop{N\!L}\nolimits}
\def\CH{\mathop{\mathrm{CH}}\nolimits}
\def\proetale{{pro\text{-}\acute{e}tale}}
\def\etale{{\acute{e}tale}}
\def\QCoh{\mathit{QCoh}}
\def\Ker{\mathop{\mathrm{Ker}}}
\def\Im{\mathop{\mathrm{Im}}}
\def\Coker{\mathop{\mathrm{Coker}}}
\def\Coim{\mathop{\mathrm{Coim}}}

% Boxtimes
%
\DeclareMathSymbol{\boxtimes}{\mathbin}{AMSa}{"02}

%
% Macros for moduli stacks/spaces
%
\def\QCohstack{\mathcal{QC}\!\mathit{oh}}
\def\Cohstack{\mathcal{C}\!\mathit{oh}}
\def\Spacesstack{\mathcal{S}\!\mathit{paces}}
\def\Quotfunctor{\mathrm{Quot}}
\def\Hilbfunctor{\mathrm{Hilb}}
\def\Curvesstack{\mathcal{C}\!\mathit{urves}}
\def\Polarizedstack{\mathcal{P}\!\mathit{olarized}}
\def\Complexesstack{\mathcal{C}\!\mathit{omplexes}}
% \Pic is the operator that assigns to X its picard group, usage \Pic(X)
% \Picardstack_{X/B} denotes the Picard stack of X over B
% \Picardfunctor_{X/B} denotes the Picard functor of X over B
\def\Pic{\mathop{\mathrm{Pic}}\nolimits}
\def\Picardstack{\mathcal{P}\!\mathit{ic}}
\def\Picardfunctor{\mathrm{Pic}}
\def\Deformationcategory{\mathcal{D}\!\mathit{ef}}
\usepackage{cleveref}
\crefname{originallemma}{original lemma}{original lemmas}
\Crefname{originallemma}{Original lemma}{Original lemmas}
\makeatletter
\newlabel{lemma-help-with-powers}{{0545}{source}{}{sourcepointer-1}{}}
\newlabel{lemma-help-with-powers@cref}{{[originallemma][0][]0545}{[1][0][]source}{}{}{}}
\makeatother
\crefname{originallemma}{original lemma}{original lemmas}
\Crefname{originallemma}{Original lemma}{Original lemmas}
\makeatletter
\newlabel{lemma-image-dense-generic-points}{{00FL}{source}{}{sourcepointer-2}{}}
\newlabel{lemma-image-dense-generic-points@cref}{{[originallemma][0][]00FL}{[1][0][]source}{}{}{}}
\makeatother
\crefname{originallemma}{original lemma}{original lemmas}
\Crefname{originallemma}{Original lemma}{Original lemmas}
\makeatletter
\newlabel{lemma-p-ring-map-field}{{0BR9}{source}{}{sourcepointer-3}{}}
\newlabel{lemma-p-ring-map-field@cref}{{[originallemma][0][]0BR9}{[1][0][]source}{}{}{}}
\makeatother
\crefname{originallemma}{original lemma}{original lemmas}
\Crefname{originallemma}{Original lemma}{Original lemmas}
\makeatletter
\newlabel{lemma-powers}{{0BR8}{source}{}{sourcepointer-4}{}}
\newlabel{lemma-powers@cref}{{[originallemma][0][]0BR8}{[1][0][]source}{}{}{}}
\makeatother
\crefname{originallemma}{original lemma}{original lemmas}
\Crefname{originallemma}{Original lemma}{Original lemmas}
\makeatletter
\newlabel{lemma-surjective-spec-radical-ideal}{{00FI}{source}{}{sourcepointer-5}{}}
\newlabel{lemma-surjective-spec-radical-ideal@cref}{{[originallemma][0][]00FI}{[1][0][]source}{}{}{}}
\makeatother
\title{Generators with p-power and p-multiple conditions plus locally nilpotent kernel give a universal homeomorphism}
\author{The Stacks Project Authors}
\date{stacks-project:algebra.tex:local-snapshot-20261012; GNU FDL 1.2 or later}
\begin{document}
\maketitle
\noindent\textit{Editorial note.} All delivered mathematical text is the author's own, copied from the pinned Stacks Project chapter file algebra.tex; no statement or proof step was rewritten, reordered or removed.\par\medskip
\noindent\textit{Editorial note.} Preamble provenance, as required for this source: the frozen snapshot stages the chapter's own preamble file beside it (preamble.tex, 217 lines, 8116 bytes, the same shared Stacks-wide preamble the published chow and etale-cohomology cards use); the first line of the pinned chapter file is \textbackslash{}input\{preamble\}, so no preamble text lives inside the chapter itself. The card takes that file's lines 8-18 (the verbatim package and the three \textbackslash{}comment-based environments reference/slogan/history at lines 9-12, section 13.2, together with the Xy-pic setup) and its lines 154 to the runtime end of the file (the native theorem declarations and mathematical macros). Not one character was changed, no package was added, no macro was defined, and the excerpt serves only as preamble, never as a span coordinate.\par\medskip
\noindent\textit{Editorial note.} The same-chapter pointers are located at their \textbackslash{}label lines inside the pinned chapter file; the cross-chapter pointers are located in the project's own tag list staged in the snapshot (sources/stacks-chow/tags.txt), at the verbatim lines carrying their canonical tags.\par\medskip
\noindent\textit{Editorial note.} The counted proof refers to same-chapter statements by label; they are recorded as navigation stubs pointed at their \textbackslash{}label lines, and the arguments using them stay inside the counted span.\par\medskip
\noindent\textit{Editorial note.} The counted spans contain no \textbackslash{}cite, so the card carries no bibliography.\par\medskip
\noindent\textit{Allowed standard prerequisites (editorial).}\begin{itemize}
\item The set of elements satisfying the p-power conditions. Introduced in the counted proof at lines 11101-11104. Isolates the elements whose images in residue fields are controlled, which is the set shown to be everything.
\item Local nilpotence of the kernel. Used at the start of the counted proof (lines 11100-11101). Allows elements of the kernel to be discarded, which is what keeps the generated set closed under the operations.
\end{itemize}\medskip
\ifdefined\hypertarget\hypertarget{sourcepointer-1}{}\fi
\noindent\textit{Source reference.} Original lemma 0545: sources/stacks-algebra/pinned-original-source.tex, source lines 11031--11031. Navigation pointer for \textbackslash{}ref\{lemma-help-with-powers\} used inside the counted span (line 11112); the argument that consumes it is retained in the span.\par\medskip
\ifdefined\hypertarget\hypertarget{sourcepointer-2}{}\fi
\noindent\textit{Source reference.} Original lemma 00FL: sources/stacks-algebra/pinned-original-source.tex, source lines 6138--6138. Navigation pointer for \textbackslash{}ref\{lemma-image-dense-generic-points\} used inside the counted span (line 11121); the argument that consumes it is retained in the span.\par\medskip
\ifdefined\hypertarget\hypertarget{sourcepointer-3}{}\fi
\noindent\textit{Source reference.} Original lemma 0BR9: sources/stacks-algebra/pinned-original-source.tex, source lines 11065--11065. Navigation pointer for \textbackslash{}ref\{lemma-p-ring-map-field\} used inside the counted span (line 11095); the argument that consumes it is retained in the span.\par\medskip
\ifdefined\hypertarget\hypertarget{sourcepointer-4}{}\fi
\noindent\textit{Source reference.} Original lemma 0BR8: sources/stacks-algebra/pinned-original-source.tex, source lines 10935--10935. Navigation pointer for \textbackslash{}ref\{lemma-powers\} used inside the counted span (line 11101); the argument that consumes it is retained in the span.\par\medskip
\ifdefined\hypertarget\hypertarget{sourcepointer-5}{}\fi
\noindent\textit{Source reference.} Original lemma 00FI: sources/stacks-algebra/pinned-original-source.tex, source lines 6039--6039. Navigation pointer for \textbackslash{}ref\{lemma-surjective-spec-radical-ideal\} used inside the counted span (line 11118); the argument that consumes it is retained in the span.\par\medskip
\begin{lemma}
\label{lemma-p-ring-map}
Let $\varphi : R \to S$ be a ring map. Let $p$ be a prime number. Assume
\begin{enumerate}
\item[(a)] $S$ is generated as an $R$-algebra by elements $x$ such
that there exists an $n > 0$ with $x^{p^n} \in \varphi(R)$ and
$p^nx \in \varphi(R)$, and
\item[(b)] $\Ker(\varphi)$ is locally nilpotent,
\end{enumerate}
Then $\varphi$ induces a homeomorphism of spectra and induces
residue field extensions satisfying the equivalent conditions
of Lemma \ref{lemma-p-ring-map-field}. For any ring map $R \to R'$
the ring map $R' \to R' \otimes_R S$ also satisfies (a) and (b).
\end{lemma}

\begin{proof}
Assume (a) and (b). Note that (b) is equivalent to condition (2)
of Lemma \ref{lemma-powers}. Let $T \subset S$ be the set of
elements $x \in S$ such that there exists an
integer $n > 0$ such that $x^{p^n} , p^n x \in \varphi(R)$.
We claim that $T = S$. This will prove that condition (1) of
Lemma \ref{lemma-powers} holds and hence $\varphi$ induces
a homeomorphism on spectra.
By assumption (a) it suffices to show that $T \subset S$ is an $R$-sub algebra.
If $x \in T$ and $y \in R$, then it is clear that $yx \in T$.
Suppose $x, y \in T$ and $n, m > 0$ such that
$x^{p^n}, y^{p^m}, p^n x, p^m y \in \varphi(R)$.
Then $(xy)^{p^{n + m}}, p^{n + m}xy \in \varphi(R)$
hence $xy \in T$. We have $x + y \in T$ by Lemma \ref{lemma-help-with-powers}
and the claim is proved.

\medskip\noindent
Since $\varphi$ induces a homeomorphism on spectra, it is in particular
surjective on spectra which is a property preserved under any base change, see
Lemma \ref{lemma-surjective-spec-radical-ideal}.
Therefore for any $R \to R'$ the kernel of the ring map
$R' \to R' \otimes_R S$ consists of nilpotent elements, see
Lemma \ref{lemma-image-dense-generic-points},
in other words (b) holds for $R' \to R' \otimes_R S$.
It is clear that (a) is preserved under base change.
Finally, the condition on residue fields follows from (a)
as generators for $S$ as an $R$-algebra map to generators for
the residue field extensions.
\end{proof}

\end{document}
